% Transcription — fonds Alexandre Grothendieck, shelfmark 101, batch 3 (pages 41-42).
% Established from the facsimile archives/batches/101/batch-03.pdf (a mirror of
% https://grothendieck.umontpellier.fr/101.pdf), per the skill
% transcribe-grothendieck. The batch file carries no cover sheet: PDF page k
% is page 40+k of the folder; the archivists' pencil numbers bottom-left
% (« 41 », « 42 ») agree. The folder has 42 pages, so this batch is its last
% two; it is short because the folder is.
%
% Pass: Opus 5.5 (claude-opus-5-5), 2026-09-24 — first pass, unchecked against the pages by a human.
%
% The folder sits in the inventory group "69-89" « Géométrie et topologie
% combinatoire », whose subject does not match this folder's (as for folder
% 100); \dating{} carries the folder's own date verbatim, and nothing is
% concluded from the group.
%
% Both pages are transcribed; nothing is skipped. No one else's text is on
% either leaf.
%
% What the batch is. It continues page 40 (batch 2): A a local ring, quotient
% of a regular ring A', n = dim A, k the residue field, an ideal of
% definition \mathfrak{N}, a system of parameters (x_1, ..., x_n), and a list
% of « Conditions » (i)-(v) on a module M (Ext^n(k,M) ≃ k, the functor
% N -> Ext^n(N,M) dualising or exact, H^n(M) dualising/injective, « A est
% C.-M. et M ≃ Ω_A »).
%   41  Blue ink, with later additions in black. A « Proposition »: if
%       N -> Ext^i(N,M) (N of finite length) is exact on one side, then
%       Ext^i(N,M) ≃ Hom(N, H^i(M)); if it is exact, it is either zero or
%       i = n with \widehat{H^n(M)} ≃ A^m, whence Ext^n(k,M) ≃ k^m and
%       Ext^i(k,M) = 0 for i > n; a black-ink question asks whether A is then
%       Cohen-Macaulay with M ≃ Ω_A^m. Below a rule, the same with the functor
%       Ext^{i-1}(-,M) right exact: either Ext^{i-1}(k,M) = 0 or
%       Ext^i(k,M) = 0 ⟺ H^i(M) = 0; a marginal note drops the finiteness of
%       M. The page ends on 0 -> M -> J -> M' -> 0 and nothing more.
%   42  Grey-black ink throughout. The conditions of page 40 rewritten for
%       M = A: (i), a struck (i bis), (ii), (iii), (iii bis), (iv),
%       (iv bis), (v), (v bis), (v ter), (vi) A/Σ x_i A autodual,
%       (vii) \bar A ≃ \bar Ω_A; a boxed programme; a scheme of implications
%       between them (drawn as a tikzcd); « A est C-M. et \bar Ω_A = \bar A »;
%       then the start of a proof that (vii) implies Cohen-Macaulay, by
%       induction on n through a non-zero-divisor x, via the spectral
%       sequence Ext^i_A(A/x, Ext^j_{A'}(A,A')) => Ext^*_{A'}(A/x,A'), ending
%       on 0 -> Ext^1_A(A/x, Ω_A) -> Ω_{A/x} -> Hom_A(A/x, Ω^{(1)}_A) -> 0.
%       The folder stops there.
%
% Notation fixed for the batch. Ext and Hom without subscript where he writes
% none; H^i(M) as he writes it (local cohomology, never indexed by the ideal
% on these pages); his barred A and Ω as \overline{A}, \overline{\Omega}_A
% (on page 42 the bar over Ω is sometimes wavy; not distinguished); the hat
% over H^n(M), \hat A, A/\mathfrak{N} as \widehat{}; the ideal of definition
% as \mathfrak{N}, as batch 2 reads the same glyph on pages 35 and 40; his Ω^{(j)}_A as written. The abbreviations « C.-M. »,
% « C-M. », « Coh. Mac. », « équidim. », « long. finie » are kept.
%
% Cross-references. Page 42's labels (i)-(vii) are those of its own list,
% which re-runs page 40's for M = A; no explicit page reference.
\documentclass[11pt,a4paper]{article}
% Le préambule commun à toutes les transcriptions du fonds Grothendieck.
%
% Il tient en une idée : **l'incertitude se compose, elle ne se résout pas.**
% Une transcription qui lisse un mot douteux a détruit la seule chose pour
% laquelle on la faisait — un lecteur doit pouvoir savoir, à chaque endroit,
% ce qui était sur le feuillet et ce qui a été ajouté. Les six macros qui
% suivent sont donc l'appareil critique, et non de la décoration.
%
% Les mêmes commandes sont reconnues par `scripts/render.mjs`, qui produit la
% vue de lecture du site. Une macro ajoutée ici doit l'être aussi là-bas,
% faute de quoi le rendu échouera bruyamment — ce qui est voulu : un rendu
% silencieusement faux est pire qu'un rendu absent.

\ProvidesPackage{grothendieck}[2026/08/08 Transcriptions du fonds Grothendieck]

\usepackage{amsmath,amssymb,amsthm}
\usepackage{mathrsfs}
\usepackage{stmaryrd}
% Les diagrammes commutatifs sont partout dans ce fonds : c'est la langue
% même de la géométrie algébrique de Grothendieck, et une transcription qui
% les rendrait en prose ne transcrirait rien.
\usepackage{tikz-cd}
% Français par défaut — c'est la langue des feuillets — mais l'anglais est
% chargé aussi : la traduction et le résumé sont des documents anglais, et la
% typographie française (espaces avant les deux-points) y serait une faute.
% Ces documents basculent par \selectlanguage{english} en tête de corps.
\usepackage[english,french]{babel}
\usepackage{fontspec}
\usepackage{geometry}
\geometry{a4paper,margin=25mm,marginparwidth=28mm}
\usepackage{marginnote}
\usepackage{xcolor}
% ulem, not soul. `soul`'s \st and \ul are built on a character-by-character
% scan that XeLaTeX's font machinery breaks: the document fails with an
% undefined control sequence rather than degrading. ulem does the same two jobs
% and is Unicode-engine safe. `normalem` stops it hijacking \emph, which the
% transcriptions use for Grothendieck's own emphasis.
\usepackage[normalem]{ulem}
\usepackage{hyperref}
\hypersetup{colorlinks=true,linkcolor=black,urlcolor=black,pdfborder={0 0 0}}

% --- Métadonnées du lot -----------------------------------------------------
% Reprises telles quelles par le script de rendu, qui les affiche en tête de
% la vue de lecture. Elles fixent ce que le fichier prétend couvrir : sans
% elles, une transcription flotte sans point d'attache dans un fonds de seize
% mille feuillets.

\newcommand{\folder}[1]{\gdef\grothFolder{#1}}
\newcommand{\batch}[1]{\gdef\grothBatch{#1}}
\newcommand{\leaves}[2]{\gdef\grothFirst{#1}\gdef\grothLast{#2}}
\newcommand{\foldertitle}[1]{\gdef\grothFoldertitle{#1}}
\gdef\grothFolder{?}\gdef\grothBatch{?}\gdef\grothFirst{?}\gdef\grothLast{?}\gdef\grothFoldertitle{}

% --- Le résumé --------------------------------------------------------------
%
% Il ouvre la lecture modernisée, sous le titre « Résumé », et il est composé
% en petit. Ce n'est pas un ornement : c'est le seul passage du document qui
% ne soit pas des mathématiques, et un lecteur doit pouvoir le reconnaître
% avant de l'avoir lu — puis le sauter s'il connaît déjà le sujet, ou s'y
% arrêter s'il ne le connaît pas. Le filet qui le suit marque la reprise du
% corps.

\newenvironment{resume}
  {\small\setlength{\parskip}{0.45em}}
  {\par\medskip\hrule\medskip}

% --- Mots-clés ---------------------------------------------------------------
%
% En anglais, en clôture du résumé de la lecture modernisée : le vocabulaire
% moderne sous lequel chercher ce que les feuillets construisent. Le manifeste
% du site (`npm run manifest`) les extrait pour étiqueter la cote entière —
% cette ligne est la source unique des tags, il n'y a pas de fichier à côté.
\newcommand{\keywords}[1]{\par\smallskip\noindent
  {\footnotesize\textsc{Keywords} --- \textit{#1}}\par}

% --- L'appareil critique ----------------------------------------------------

% Le numéro de feuillet, en marge. C'est l'ancre qui permet de régler un
% désaccord sur une formule en regardant *un* feuillet plutôt qu'une chemise
% de six cents. Le site s'en sert aussi pour tourner les pages du fac-similé
% quand on fait défiler la transcription.
\newcommand{\leaf}[1]{%
  \marginnote{\footnotesize\color{gray!70}\textsf{#1}}%
  \label{leaf:#1}\ignorespaces}

% Illisible : on ne devine pas. Un mot inventé qui se lit comme les autres est
% le pire résultat possible.
\newcommand{\ill}{\textcolor{red!60!black}{[\,\ldots\,]}}

% Lecture douteuse : le mot est donné, et signalé comme incertain.
\newcommand{\uncertain}[1]{\uline{#1}}

% Ajout de l'éditeur — une lettre manquante, un mot sous-entendu.
\newcommand{\add}[1]{\textcolor{blue!50!black}{[#1]}}

% Biffé par Grothendieck lui-même. Ce qu'il a rayé fait partie du document :
% une rature dit souvent où la pensée a changé de direction.
\newcommand{\struck}[1]{\sout{#1}}

% Note du transcripteur, sur l'état matériel du feuillet.
\newcommand{\note}[1]{\par\smallskip{\small\color{gray!60!black}\textit{[#1]}}\par\smallskip}

% Note marginale de Grothendieck, distincte du corps du texte.
\newcommand{\marginal}[1]{\par\smallskip\noindent\hspace{1em}%
  {\small\color{gray!60!black}#1}\par\smallskip}

% --- Titre ------------------------------------------------------------------

\AtBeginDocument{%
  \begin{center}
    {\large\bfseries Fonds Alexandre Grothendieck --- cote n\textsuperscript{o}~\grothFolder}\\[2pt]
    {\itshape\grothFoldertitle}\\[4pt]
    {\small feuillets \grothFirst--\grothLast\ (lot \grothBatch)}\\[2pt]
    {\footnotesize\color{gray!60!black}Transcription établie d'après le fac-similé numérisé
     par l'Université de Montpellier.\\
     Les lectures incertaines sont soulignées, les illisibles notées [\,\ldots\,],
     les ajouts entre crochets.}
  \end{center}
  \vspace{1em}}

\endinput


\folder{101}
\batch{3}
\pages{41}{42}
\dating{[à partir de 1967-1968]}
\watermark{Édition de démonstration}
\foldertitle{Cohomologie locale : notes manuscrites (s.d.).}

\begin{document}

\section{[Exactitude des foncteurs Ext ; le cas du module A]}
\note{titre de l'éditeur, entre crochets : les feuillets n'en portent aucun.
Le feuillet 41 enchaîne sur les « Conditions » du feuillet 40 (lot 2) : $A$
local, quotient d'un anneau régulier $A'$, $n = \dim A$, $k$ le corps
résiduel, $M$ un $A$-module.}

\page{41}
\emph{Proposition.} Supposons le foncteur
\[
  N \mapsto \operatorname{Ext}^i(N, M) \qquad (N \text{ de longueur finie})
\]
\emph{\uncertain{exact à g.}} Alors \uncertain{on a} :
\[
  \operatorname{Ext}^i(N, M) \simeq \operatorname{Hom}(N, H^i(M))
\]
[\struck{\uncertain{a-t-on}} \struck{\uncertain{donc} \ill{}}
$\operatorname{Ext}^i(k, M)$ \ill{} ou $\operatorname{Ext}^{i-1}(k, M) = 0$ ??]
\note{crochet en marge droite ; ses deux premières lignes sont biffées, et une
flèche épaisse, repassée, les traverse vers la droite.}

Supposons le foncteur \emph{exact}, alors ou bien il est identiquement nul,
[i.e. \struck{$H^i(M) = 0$, \uncertain{ou encore}} $\operatorname{Ext}^i(k,
M) = 0$, ce qui implique \uncertain{pb.t} $i > n$ \emph{ou} $i <
\uncertain{\mathrm{codim.\ coh.}}\ M$], ou bien $i = n$, $\widehat{H^n(M)}
\simeq A^m$. Alors $M$ est \emph{\uncertain{fidèle}}, $\operatorname{Ext}^n(k,
M) \simeq k^m$. De plus, on voit que $\operatorname{Ext}^i(k, M) = 0$ pour $i
> n$ [et cela implique-t-il que $A$ est \uncertain{Coh. Mac.}, et $M \simeq
\Omega_A^m$]
\note{« pb.t » est une abréviation que l'on ne résout pas ; « codim. coh. »
est lu sur des premières lettres surchargées. Le crochet final est un ajout à
l'encre noire ; sous « Coh. Mac. », un mot à l'encre bleue est biffé
(\struck{\ill{}}) et un ajout interlinéaire, bleu, ne se lit pas.}

\marginal{$M$ non \uncertain{néc.} \uncertain{de type fini}}
Supposons le foncteur $N \mapsto \operatorname{Ext}^{\uncertain{i-1}}(N, M)$
\emph{exact} \uncertain{à} \emph{droite}. Je dis que l'on a ou bien
$\operatorname{Ext}^{\uncertain{i-1}}(k, M) = 0$, ou bien
\[
  \operatorname{Ext}^{\uncertain{i}}(k, M) = 0 \Longleftrightarrow H^i(M) = 0 .
\]
En \uncertain{effet}, \ill{} \ill{}
\[
  \operatorname{Ext}^i(N, M) \simeq \operatorname{Hom}(N, H^i(M))
\]
\note{les trois exposants de $\operatorname{Ext}$ de ce paragraphe sont
surchargés (un $i$ corrigé en $i-1$, ou l'inverse) et leur lecture est
douteuse. Le premier « ou bien » est ajouté à l'encre noire ; une accolade à
l'encre noire, à gauche, embrasse les trois lignes. La note marginale est
écrite en oblique, en haut à gauche du paragraphe.}

\[
  0 \to M \to J \to M' \to 0
\]

\page{42}
(i) $\operatorname{Ext}^n(k, A) \simeq k$

\struck{(i bis) $\operatorname{Ext}^n(A/\mathfrak{N}, A) \simeq \widehat{A}$}

(ii) $\operatorname{Ext}^{n-1}(k, A) = 0$
\note{la ligne biffée, dont l'étiquette « (i bis) » surcharge un « (ii) »,
et la ligne (ii) sont entourées ensemble d'un cadre que traverse une ligne
sinueuse.}

(iii) $\operatorname{Ext}^{n+1}(k, A) = 0$

(iii bis) $\operatorname{Ext}^n(A/\mathfrak{N}, A) \simeq
\widehat{A/\mathfrak{N}}$ \quad \struck{(\ill{} \ill{} : $\lg
\operatorname{Ext}^n(A/\mathfrak{N}, A) = \lg(A/\mathfrak{N})$)}

(iv) Le foncteur $N \mapsto \operatorname{Ext}^n(N, A)$ est dualisant

(iv bis) [idem] exact ($N$ de long. finie)
\note{pour (iv bis), un trait tient lieu de la ligne (iv) répétée.}

(v) \struck{\ill{}} $H^n(A)$ dualisant

(v bis) $H^n(A)$ injectif

(v ter) $H^n(A)$ \uncertain{comonogène} [\uncertain{resp.} $N \mapsto
\operatorname{Ext}^n(N, A)$ \uncertain{exact}] et $A$ équidim.

(vi) $A/\sum x_i A$ est autodual

(vii) $\overline{A} \simeq \overline{\Omega}_A$

\marginal{\uncertain{Prouver} \ill{} \ill{} \ill{} des \uncertain{lemmes}.
(i) (iii') (v bis) (vi) impliquent que \uncertain{l'équidimensionnalité}
\ill{} \ill{}}
\note{texte encadré, en haut à droite du feuillet ; « (iii') » est lu tel
quel. La fin du cadre ne se lit pas.}

\begin{tikzcd}[column sep=small, row sep=small, nodes={font=\scriptsize}]
{[(\mathrm{vii}) \Leftrightarrow (\mathrm{v}) \Leftrightarrow (\mathrm{v\ bis}) \Leftrightarrow (\mathrm{v\ ter})]} \arrow[d, Rightarrow] & {[(\mathrm{iv\ bis}) \Leftrightarrow (\mathrm{iv})]} \arrow[l, Rightarrow] \arrow[d, Rightarrow] \\
(\mathrm{iii}) & (\mathrm{i\ bis}) \quad (\mathrm{i}) \quad (\mathrm{vi})
\end{tikzcd}

\note{schéma d'implications, redessiné : les deux groupes sont encadrés, et
« (iii) » l'est aussi. Entre « (v) » et « (v bis) », un mot est biffé par
hachures ; entre la flèche horizontale et le cadre de droite, un pavé d'encre
couvre ce qui y était écrit.}

$A$ est C-M. et $\overline{\Omega}_A = \overline{A}$
\note{écrit à gauche, sous (vii), à la hauteur du « (iii) » encadré.}

(\ill{} \ill{} \ill{} 14 \emph{\ill{}}) (\uncertain{Récurrence} \uncertain{sur}
$n$, \uncertain{triviale} si $n = 0$)
\note{ces deux parenthèses sont ajoutées au-dessus de la ligne suivante ; une
accolade relie la seconde à « $A$ est C.-M. ».}

Prouvons que (vii) implique que $A$ est C.-M. En effet, $A$ est équidim.
puisque le module $\overline{\Omega}_A$ l'est. Soit $x \in A$ \uncertain{non}
diviseur de zéro, il \uncertain{suffit} que \uncertain{diviseur de} zéro dans
$\Omega_A$. Je dis donc que
\[
  \overline{\Omega}_{A/xA} \simeq \Omega_A / x\,\Omega_A
\]
(\uncertain{d'où} \struck{$\operatorname{Ext}^{*}_{A'}$} le résultat
\uncertain{grâce à} la récurrence). En effet,
\[
  \overline{\Omega}_{A/xA} = \struck{H^{\ill{}}(A/xA) = \ill{}}
  \operatorname{Ext}^{\uncertain{n-1}}_{A'}(A/x, A')
\]
\uncertain{or}
\[
  \struck{H^{\ill{}}(A/xA) = \ill{}}
\]
\[
  \operatorname{Ext}^{*}_{A'}(A/x, A') \Longleftarrow
  \operatorname{Ext}^i_A\bigl(A/x, \underbrace{\operatorname{Ext}^{\uncertain{j}}_{A'}(A, A')}_{\Omega^{(j)}_A}\bigr)
  \qquad \ill{}
\]
\note{l'exposant du premier $\operatorname{Ext}_{A'}$ pourrait se lire
$r-n$ ; au-dessus de la ligne biffée est écrit « $r-n+1$ », et au-dessus de
$\operatorname{Ext}^{j}_{A'}(A, A')$ est écrit « $r-n$ ». Dans la première
égalité, une flèche courbe renvoie de la partie biffée à
$\operatorname{Ext}_{A'}$.}

\uncertain{d'où}
\[
  \Omega^{(*)}(A/x) \Longleftarrow \operatorname{Ext}^i_A(A/x, \Omega^{(j)}_A)
\]
[\uncertain{Bien} \ill{} \ill{} par $A/x$ \ill{} \ill{} \ill{} $A$-module
\uncertain{quelconque}]
\marginal{\emph{Lemme} \ill{} \ill{}}
\note{le crochet est en marge droite, sur quatre lignes ; la note marginale,
à gauche, est soulignée deux fois pour « Lemme », le reste ne se lit pas.}

d'où
\[
  0 \to \operatorname{Ext}^1_A(A/x, \Omega_A) \to \Omega_{A/x} \to
  \operatorname{Hom}_A(A/x, \Omega^{(1)}_A) \to 0 .
\]
\note{$\Omega_{A/x}$ est souligné d'un trait incurvé ; un signe, peut-être
biffé, surmonte $\operatorname{Hom}_A$. Le dossier s'arrête sur cette ligne.}

\end{document}
